\documentclass[11pt]{article}
\usepackage[margin=1in]{geometry}
\usepackage{amsmath}
\usepackage{graphicx}
\usepackage{booktabs}
\usepackage{listings}
\usepackage{xcolor}
\usepackage{hyperref}
\usepackage{enumitem}

\lstset{
  basicstyle=\ttfamily\small,
  breaklines=true,
  columns=fullflexible,
  keywordstyle=\color{blue},
  commentstyle=\color{gray},
  stringstyle=\color{teal},
  frame=single,
  numbers=left,
  numberstyle=\tiny\color{gray},
  xleftmargin=1.5em,
}

\title{Microstrip Ring Resonator on Polyimide:\\
Model Development and Modification Guide}
\author{openEMS Ring Resonator Dk/Loss Extraction Model}
\date{\today}

\begin{document}
\maketitle
\tableofcontents

\section{Purpose}

This document describes the openEMS/CSXCAD full-wave FDTD model implemented
in \texttt{07\_ring\_resonator\_polyimide\_dk\_loss.py}. The model simulates a
microstrip ring resonator, weakly gap-coupled to a single straight through
feed line, fabricated on a thin polyimide substrate. A single wideband
Gaussian excitation spanning 100~MHz--67~GHz is used so that every harmonic
resonance of the ring appears as a notch in the transmitted S-parameter
\(|S_{21}|\). From the measured resonant frequencies, $-3$~dB bandwidths, and
insertion losses, the script extracts the substrate's relative permittivity
(Dk) and loss tangent ($\tan\delta$) as a function of frequency --- the
classic \emph{ring-resonator method} used for microwave/mm-wave laminate
characterization.

This is a practical, general-purpose starting point: the same script
structure can be adapted to other substrates, stack-ups, and ring/feed
geometries by editing a small number of parameters near the top of the file,
as described in Section~\ref{sec:modifying}.

\section{Model Development History}

The model went through several rounds of debugging before it converged on
physically correct S-parameters. These are documented here because the
underlying causes are generic FDTD/microstrip modeling pitfalls that are
easy to reintroduce when the geometry or substrate stack is modified --
Section~\ref{sec:modifying} calls out where each rule must be respected
again.

\begin{enumerate}[leftmargin=*]

\item \textbf{Absorbing boundary choice.} The \texttt{MUR} boundary
condition was found to be numerically unstable (unbounded exponential
energy growth) for this broadband (100~MHz--67~GHz) case, especially in
combination with a PEC ground plane sitting exactly on the \texttt{zmin}
domain boundary. The fix was to use \texttt{PML\_8} on all open boundaries
and \texttt{PEC} only on the \texttt{zmin} face (where the ground plane
already terminates the field):
\begin{lstlisting}[language=Python]
FDTD.SetBoundaryCond(["PML_8","PML_8","PML_8","PML_8","PEC","PML_8"])
\end{lstlisting}

\item \textbf{Port margin.} A lumped port's current probe measures a closed
H-field loop that needs a valid FDTD cell on \emph{both} sides of the probe.
A port placed exactly on the outermost mesh line of the simulation domain
has no cell on its outward side, so the measured current is pure numerical
noise, corrupting $S_{11}$/$S_{21}$ even though the simulation itself
remains stable. The substrate and ground plane are extended by a
\texttt{margin} (1~mm) beyond each port so the ports sit a few cells inside
the domain.

\item \textbf{Feed line must stop exactly at the port planes.} The
metal trace must \emph{not} extend into the margin region past $x=0$ or
$x=\texttt{sub\_l}$: a continuous metal sheet running straight over a lumped
port's vertical probe shorts out the port's series resistor, again causing
near-total, frequency-independent reflection.

\item \textbf{Open-boundary air gap (z-direction).} Without sufficient air
space above the trace, the PML/\texttt{zmax} boundary sits right on (or a
fraction of a millimeter above) the metal, which either silently falls back
to PEC behavior (too few $z$-lines for a full 8-cell PML) or absorbs the
near field before it can couple properly. The fix grades the $z$-mesh
manually above the trace with a geometrically growing sequence (ratio 1.3)
out to \texttt{z\_air = max(10h, 5w, 1.5)}~mm of clearance.

\item \textbf{Lateral clearance (y-direction).} The same principle applies
in-plane: the feed line's outer edge must be kept well clear of the
\texttt{PML\_8} boundary at $y=0$. A microstrip's fringing E-fields extend
several substrate-heights beyond the trace edge; truncating that too close
to an absorbing boundary parasitically loads the line (acts like a nearby
ground wall), severely distorting its characteristic impedance and causing
large, frequency-independent reflection at \emph{every} frequency --
including very low ones where a matched line should behave almost ideally.
The fix introduces
\begin{lstlisting}[language=Python]
y_clear = max(10 * h, 5 * w, 0.5)   # [mm] lateral clearance to y=0 boundary
\end{lstlisting}
and uses it to size the substrate width and feed-line offset.

\item \textbf{Mesh-line deduplication vs.\ exact port boundaries.} A
deduplication pass (\texttt{dedupe\_lines}) merges near-coincident mesh
lines to avoid tiny sliver cells that would otherwise force an unstable CFL
time step. This can silently merge an exact, manually inserted port-boundary
line into a nearby fine-grid point, drifting the port's edge away from the
intended trace edge. The fix re-inserts the exact port-boundary lines
(\texttt{mesh.AddLine("x", [0, sub\_l])} and the $y$ edges of the feed line)
\emph{after} the dedup pass.

\item \textbf{Kappa reference frequency.} openEMS's \texttt{kappa} material
parameter is a plain (frequency-independent) Ohmic conductivity, so the
\emph{effective} loss tangent it produces scales as
$\tan\delta_{\mathrm{eff}}(f) = \tan\delta_{\mathrm{nominal}} \cdot
f_0/f$ -- it is \textbf{not} constant with frequency. Using the excitation's
arithmetic-mean center frequency as $f_0$ badly overestimates loss at the
low end of a log-wide band. The fix uses the \emph{geometric} mean of the
band edges instead:
\begin{lstlisting}[language=Python]
f0_kappa = float(np.sqrt(f_start * f_stop))
kappa = tanD_nominal * 2 * np.pi * f0_kappa * er_nominal * eps0
\end{lstlisting}
This minimizes the worst-case log-frequency error across the band but does
not eliminate it -- see Section~\ref{sec:limitations}.

\item \textbf{Resonance signature: dips, not peaks.} For this single-tap,
gap-coupled ring topology, each ring resonance couples energy \emph{out of}
the through line and into the ring, producing a notch (local minimum) in
$|S_{21}|$ at resonance, not a peak. The post-processing peak finder
searches for local minima, not maxima.

\item \textbf{Harmonic-index assignment.} The lowest-order ring harmonics
are the most weakly coupled and can be too shallow to cross the detection
threshold, especially before the simulation has fully converged. If the
detected notches were simply numbered $n=1,2,3,\dots$ in ascending frequency
order, a missing low-order harmonic would silently mislabel every harmonic
above it. The script instead estimates $n$ directly from the physical
resonance condition (Eq.~\ref{eq:resonance}) using the nominal effective
permittivity as an initial guess, then rounds to the nearest integer.

\end{enumerate}

\section{Model Structure}

\subsection{Geometry overview}

\begin{center}
\begin{tabular}{@{}ll@{}}
\toprule
\textbf{Element} & \textbf{Description} \\
\midrule
Ground plane & Full-size PEC sheet at $z=0$, covering the substrate footprint (plus margin) \\
Substrate & Single dielectric slab, thickness $h$, from $z=0$ to $z=h$ \\
Feed line & Straight microstrip trace of width $w$, thickness $t$, at $z \in [h, h+t]$ \\
Ring & Cylindrical-shell (annulus) trace, mean radius $r_{\mathrm{mean}}$, width $w$, same $z$ range \\
Ports & Two lumped 50~$\Omega$ ports, one at each end of the feed line ($z$-directed) \\
\bottomrule
\end{tabular}
\end{center}

The ring is placed alongside the feed line, separated by a coupling gap
\texttt{gap}, so that the ring's outer edge is a distance \texttt{gap} from
the feed line's edge (edge-coupled, not tap-coupled). Only a single feed
line couples to the ring (a ``one-port-pair, single-tap'' topology): energy
is launched at port~1, mostly transmitted straight through to port~2, with a
small fraction resonantly coupled into the ring at each of its harmonic
frequencies.

\subsection{Key parameters in the script}

\begin{center}
\begin{tabular}{@{}lll@{}}
\toprule
\textbf{Variable} & \textbf{Meaning} & \textbf{Default} \\
\midrule
\texttt{f\_start}, \texttt{f\_stop} & Simulation frequency band & 0.1--67~GHz \\
\texttt{er\_nominal} & Substrate relative permittivity (Dk) & 3.5 \\
\texttt{tanD\_nominal} & Substrate loss tangent & 0.008 \\
\texttt{h} & Substrate thickness [mm] & 0.1 \\
\texttt{w} & Trace width [mm], synthesized for 50~$\Omega$ & (computed) \\
\texttt{trace\_t} & Copper trace thickness [mm] & 0.017 (0.5~oz) \\
\texttt{f1\_target} & Target fundamental ring resonance & 8.5~GHz \\
\texttt{r\_mean} & Ring mean radius [mm] & (computed from \texttt{f1\_target}) \\
\texttt{gap} & Ring-to-feed-line coupling gap [mm] & 0.12 \\
\texttt{feed\_len} & Feed stub length each side [mm] & 4.0 \\
\texttt{margin} & Port/substrate margin past the ports [mm] & 1.0 \\
\texttt{y\_clear} & Lateral clearance, trace to PML [mm] & $\max(10h, 5w, 0.5)$ \\
\texttt{z\_air} & Vertical clearance, trace to PML [mm] & $\max(10h, 5w, 1.5)$ \\
\bottomrule
\end{tabular}
\end{center}

\subsection{Trace width synthesis}

The feed-line/ring trace width $w$ is synthesized for a 50~$\Omega$
characteristic impedance using the closed-form Hammerstad synthesis formula,
given the substrate's nominal $\varepsilon_r$ and thickness $h$:
\[
A = \frac{Z_0}{60}\sqrt{\frac{\varepsilon_r+1}{2}}
   + \frac{\varepsilon_r - 1}{\varepsilon_r + 1}
     \left(0.23 + \frac{0.11}{\varepsilon_r}\right), \qquad
\frac{W}{h} = \frac{8 e^{A}}{e^{2A} - 2} \quad (W/h < 2).
\]
This is only valid for the narrow-trace ($W/h < 2$) regime; for thicker
substrates or higher-Dk materials where the synthesized $W/h \geq 2$, the
wide-trace Hammerstad formula should be substituted instead (see
Section~\ref{sec:modifying}).

\subsection{Ring sizing and the resonance condition}
\label{sec:resonance}

The ring's mean radius is chosen so its fundamental ($n=1$) resonance sits
near a target frequency \texttt{f1\_target} (default 8.5~GHz), spreading
roughly seven harmonics across the 100~MHz--67~GHz simulation band. The
governing relation is the ring's circumferential standing-wave condition:
\begin{equation}
n\,\lambda_g(f_n) = 2\pi r_{\mathrm{mean}}
\quad\Longrightarrow\quad
\varepsilon_{\mathrm{eff}, n} =
\left(\frac{n\,c_0}{2\pi r_{\mathrm{mean}} f_n}\right)^{2},
\label{eq:resonance}
\end{equation}
where $c_0$ is the speed of light in vacuum and $f_n$ is the $n$-th measured
resonant frequency. This same relation, inverted, is used both to size the
ring initially (given a nominal $\varepsilon_{\mathrm{eff}}$) and to
recover $\varepsilon_{\mathrm{eff}, n}$ from each measured resonance.

\subsection{Dk and loss-tangent extraction}

For each detected $|S_{21}|$ notch:
\begin{enumerate}
\item Determine $f_n$ and its harmonic index $n$ (Eq.~\ref{eq:resonance}).
\item Compute $\varepsilon_{\mathrm{eff}, n}$ from Eq.~\ref{eq:resonance}.
\item Invert the Hammerstad--Jensen effective-permittivity relation for the
known (fixed) trace $W/h$ to recover $\varepsilon_{r,n}$ (Dk):
\[
\varepsilon_{\mathrm{eff}} = \frac{\varepsilon_r + 1}{2}
  + \frac{\varepsilon_r - 1}{2} F(W/h), \qquad
  F(W/h) = \left(1 + \frac{12}{W/h}\right)^{-1/2}.
\]
\item Measure the loaded $Q_L = f_n / \mathrm{BW}_{3\mathrm{dB}}$ from the
notch's $-3$~dB bandwidth, then the unloaded $Q_u$ via the standard
insertion-loss correction $Q_u = Q_L / (1 - 10^{-\mathrm{IL_{dB}}/20})$.
\item Compute an analytic conductor-loss-only $Q_c$ from the Wheeler
incremental-inductance rule and copper surface resistance, and remove it
from $Q_u$ to isolate the dielectric-only $Q_d$:
$1/Q_d = 1/Q_u - 1/Q_c$ (only when $Q_c > Q_u$; otherwise $Q_u$ itself is
reported as an upper bound on $Q_d$, see Section~\ref{sec:limitations}).
\item Recover the loss tangent via the dielectric filling factor
$q = (\varepsilon_{\mathrm{eff}} - 1)/(\varepsilon_r - 1)$:
\[
\tan\delta_n = \frac{1}{Q_{d,n}\, q_n}.
\]
\end{enumerate}

\section{Limitations and Interpretation Notes}
\label{sec:limitations}

\begin{itemize}[leftmargin=*]
\item \textbf{Convergence}: The default \texttt{NrTS = 200000} /
\texttt{EndCriteria = 1e-5} ($-50$~dB) budget may not be reached for a
genuinely high-$Q$ ring before the timestep budget is exhausted (openEMS
will print a warning and stop at \texttt{NrTS}). An incompletely decayed
time-domain signal narrows the apparent resonance linewidth (spectral
leakage), which \emph{inflates} the apparent $Q_u$ and therefore
\emph{underestimates} the extracted $\tan\delta$. If the extracted Dk looks
reasonable but $\tan\delta$ looks too low (or increases smoothly as
frequency decreases across the band), increase \texttt{NrTS} and/or relax
\texttt{EndCriteria} slightly and re-run.
\item \textbf{Frequency-independent kappa}: because openEMS models a
constant Ohmic conductivity rather than a true frequency-independent
$\tan\delta$, the simulated substrate loss is only exactly correct at
$f = f_0\_\mathrm{kappa}$ (the geometric mean of the band) and drifts
(over- or under-representing the true material loss) away from it. This is
a modeling compromise, not a bug; it does not need to be "fixed" for
Dk extraction (which does not depend on kappa), but should be kept in mind
when interpreting the extracted $\tan\delta$ trend vs.\ frequency.
\item \textbf{Analytic conductor-loss model}: the built-in $Q_c$ estimate
is a simplified microstrip attenuation formula, not a full analytic model
(it ignores edge current crowding, surface roughness, etc.). If $Q_c$ comes
out lower than the measured $Q_u$ (physically impossible for a clean
loss-mechanism separation), the script falls back to reporting $Q_u$
directly, which is a conservative (upper-bound-on-Dk-loss,
lower-bound-on-actual-$\tan\delta$) estimate rather than a hard failure.
\end{itemize}

\section{Running the Model}

\begin{lstlisting}[language=bash]
conda activate openems
python 07_ring_resonator_polyimide_dk_loss.py
\end{lstlisting}

Simulation results (port time/frequency-domain data, and the final
S-parameter/Dk/loss plot \texttt{ring\_resonator\_dk\_loss.png}) are written
to \texttt{sim\_results/ring\_resonator/}. A full run at the default mesh
density takes on the order of a few hours of wall-clock time on a modern
multi-core workstation (roughly 2.8~million FDTD cells, tens of thousands
of timesteps).

\section{Modifying the Model}
\label{sec:modifying}

All of the changes below only require editing parameters/geometry near the
top of \texttt{07\_ring\_resonator\_polyimide\_dk\_loss.py}; the mesh, port,
and boundary-condition logic further down the file automatically adapts
(re-derives clearances, resolutions, etc.) from these parameters and does
not need to be touched for routine changes described here.

\subsection{Changing the dielectric properties}

To model a different substrate material, edit only:
\begin{lstlisting}[language=Python]
er_nominal   = 3.5     # relative permittivity (Dk)
tanD_nominal = 0.008   # loss tangent
h            = 0.1     # substrate thickness [mm]
\end{lstlisting}
Everything else (trace width synthesis, ring sizing, mesh resolution,
kappa) is automatically recomputed from these three values. If the new
material's synthesized trace width comes out with $W/h \geq 2$ (thick,
low-Dk substrates, or very thin high-Dk ones), replace the narrow-trace
Hammerstad formula (Section~\ref{sec:modifying}, trace-width synthesis)
with the wide-trace variant:
\[
\frac{W}{h} = \frac{2}{\pi}\Big[B - 1 - \ln(2B-1)
  + \frac{\varepsilon_r - 1}{2\varepsilon_r}
    \Big(\ln(B-1) + 0.39 - \frac{0.61}{\varepsilon_r}\Big)\Big],
\qquad B = \frac{377\pi}{2 Z_0 \sqrt{\varepsilon_r}}.
\]

\subsection{Changing the ring/feed geometry}

\begin{itemize}[leftmargin=*]
\item \textbf{Fundamental resonance frequency}: change \texttt{f1\_target}
(default 8.5~GHz). The ring radius \texttt{r\_mean} is automatically
recomputed from Eq.~\ref{eq:resonance} using the nominal
$\varepsilon_{\mathrm{eff}}$, so the harmonic comb shifts accordingly. Keep
in mind the number of harmonics visible in the fixed 100~MHz--67~GHz band
scales roughly as $67\,\mathrm{GHz}/f_1$.
\item \textbf{Coupling strength}: change \texttt{gap} (default 0.12~mm).
Smaller gaps couple more strongly (deeper, wider notches, lower loaded
$Q_L$); larger gaps couple more weakly (shallower notches that may fall
below the detection threshold, particularly at low frequency -- see
Section~\ref{sec:limitations}).
\item \textbf{Feed line length}: change \texttt{feed\_len} (default 4.0~mm
each side). This mainly affects total simulation domain size/run time and
has negligible effect on the extracted results as long as it stays large
enough that the port's near-field has settled before reaching the ring.
\item \textbf{Trace width / impedance}: the trace width \texttt{w} is
automatically synthesized for 50~$\Omega$ from \texttt{er\_nominal} and
\texttt{h}. To target a different characteristic impedance, change
\texttt{Z0} (default 50.0) before the synthesis block.
\item \textbf{Copper thickness}: change \texttt{trace\_t} (default
0.017~mm, i.e. 0.5~oz copper) to model a different copper weight.
\end{itemize}

Whenever the ring radius, gap, or trace width changes substantially, sanity
check that the automatically-derived clearances
(\texttt{y\_clear}, \texttt{z\_air}, \texttt{margin}) are still being applied
correctly (Section~\ref{sec:modifying} and history items 4 and 5): these are
all defined as functions of \texttt{h} and \texttt{w}, so they should scale
appropriately on their own, but always re-inspect the printed geometry
summary at the start of a run to confirm reasonable values before committing
to a multi-hour simulation.

\subsection{Adding more substrate layers}

The current model uses a single dielectric slab between the ground plane
and the trace. To model a multilayer stack-up (e.g., a core + prepreg, or a
polyimide film laminated to a carrier), replace the single
\texttt{CSX.AddMaterial("Polyimide", ...)} block with one material and one
\texttt{AddBox} call per layer, stacked in $z$:

\begin{lstlisting}[language=Python]
# Example: two-layer substrate (bottom layer + top layer under the trace)
h1, er1, tanD1 = 0.05, 3.5, 0.008   # bottom layer (polyimide)
h2, er2, tanD2 = 0.05, 4.2, 0.015   # top layer (e.g. bond-ply/adhesive)
h = h1 + h2                         # total substrate thickness used everywhere else

kappa1 = tanD1 * 2 * np.pi * f0_kappa * er1 * eps0
kappa2 = tanD2 * 2 * np.pi * f0_kappa * er2 * eps0

sub1 = CSX.AddMaterial("Polyimide_bottom", epsilon=er1, kappa=kappa1)
sub1.AddBox(priority=0, start=[-margin, 0, 0], stop=[sub_l + margin, sub_w, h1])

sub2 = CSX.AddMaterial("Polyimide_top", epsilon=er2, kappa=kappa2)
sub2.AddBox(priority=0, start=[-margin, 0, h1], stop=[sub_l + margin, sub_w, h1 + h2])

# Ground plane, feed line, ring, and ports are unaffected: they already
# reference z=0 (ground) and z=h..h+trace_t (trace) generically, so as
# long as `h` above is redefined as the TOTAL stack thickness, no other
# geometry code needs to change.
\end{lstlisting}

Important considerations when adding layers:
\begin{itemize}[leftmargin=*]
\item \textbf{Effective permittivity synthesis} (trace width, ring sizing,
$F(W/h)$ inversion) assumes a single homogeneous substrate. For a true
multilayer stack, either (a) use a weighted/parallel-plate approximate
$\varepsilon_r$ for the synthesis step only (reasonable when one layer
dominates the field energy, e.g. a thin bond-ply under a much thicker
core), or (b) treat \texttt{er\_nominal} in the synthesis formulas as an
initial guess and rely on the FDTD-measured $\varepsilon_{\mathrm{eff}, n}$
per layer stack as the authoritative result -- the resonance-based
extraction in Section~\ref{sec:resonance} works correctly regardless of how
many dielectric layers are present, since it only depends on the
\emph{measured} resonant frequencies, not on the analytic synthesis
formulas used to size the geometry initially.
\item \textbf{Mesh}: add a $z$-mesh line at every layer interface
(\texttt{mesh.AddLine("z", [0, h1, h1+h2, ...])}) so each layer boundary is
resolved exactly; the existing fine-to-coarse grading logic for the
open-boundary air gap above the trace does not need to change, since it
already starts counting from \texttt{h + trace\_t}, the top of the (now
multi-layer) stack.
\item \textbf{Reporting}: if you want per-layer Dk/loss reported separately
rather than a single lumped ``Dk'' number, you will need a more detailed
post-processing model (e.g. conformal-mapping partial-capacitance methods)
since the ring-resonator method as implemented here only recovers a single
effective $\varepsilon_r$ per resonance, consistent with a single equivalent
homogeneous substrate.
\end{itemize}

\section{File Summary}

\begin{center}
\begin{tabular}{@{}ll@{}}
\toprule
\textbf{File} & \textbf{Description} \\
\midrule
\texttt{07\_ring\_resonator\_polyimide\_dk\_loss.py} & Full model: geometry, mesh, simulation, post-processing \\
\texttt{sim\_results/ring\_resonator/} & openEMS output directory (port data, plot) \\
\texttt{ring\_resonator\_dk\_loss.png} & Final S-parameter / extracted Dk / extracted $\tan\delta$ plot \\
\bottomrule
\end{tabular}
\end{center}

\end{document}
